\section{Actual local maps and endpoint operators}
\label{sec:peps_endpoint_local_operators}

\begin{theorem}[Insertion of the recorded private input]
    \label{thm:peps_recorded_private_input}
    \lean{TNLean.PEPS.PairEffect.SourceCircuit.eval_sourcePrivatePreparation_tmul,
      TNLean.PEPS.PairEffect.SourceCircuit.eval_partialSourceWord_tmul}
    \leanok
    For a chronological branch and a recorded party product input $\psi$,
    preparing the private input while leaving the source registers unchanged
    sends a source vector $y$ to $y\otimes\psi$. Composing this preparation
    with the chronological residual gives its action on that vector.
\end{theorem}

\begin{proof}\leanok
    Each private input is prepared at its recorded register and owner.
    Framing these preparations by the source registers acts as the identity
    on the source vector. Apply the definition of composition.
\end{proof}

\begin{definition}[Local branch and physical matrices]
    \label{def:peps_actual_local_branch_matrix}
    \lean{TNLean.PEPS.PairEffect.SourceCircuit.partialSourceLocalMatrix,
      TNLean.PEPS.PairEffect.SourceCircuit.partialSourceLocalPhysicalMatrix}
    \leanok
    In chosen orthonormal bases, let $L_\eta$ be the matrix of the actual
    party-local source-to-output map at incident branch choice $\eta$,
    after the recorded private input is prepared. For a source operator $X$
    and two choices $\eta,\zeta$, the local physical matrix is
    $\operatorname{Tr}_{\rm private}(L_\eta X L_\zeta^*)$.
    The source operator need not be positive.
\end{definition}

\begin{theorem}[Agreement with the chronological restriction]
    \label{thm:peps_actual_local_branch_matrix}
    \lean{TNLean.PEPS.PairEffect.SourceCircuit.partialSourceLocalMatrix_eq_actual_restriction,
      TNLean.PEPS.PairEffect.SourceCircuit.choicesOfOriginalGateLabelsAt_eq_projectAllChoices}
    \leanok
    \uses{def:peps_actual_local_branch_matrix}
    Projecting an actual global branch choice to its incident original
    gate occurrences gives precisely the local choices used by the
    restricted chronological map. The local matrix above is its matrix
    in the chosen bases.
\end{theorem}

\begin{proof}\leanok
    The correspondence between touched occurrences and incident labels
    commutes with projection. Apply the equality of the local map and
    the actual chronological restriction to each input basis vector.
\end{proof}

\begin{definition}[Endpoint operator word]
    \label{def:peps_actual_endpoint_word}
    \lean{TNLean.PEPS.PairEffect.Word.localEndpoint,
      TNLean.PEPS.PairEffect.SourceInventory.endpointWord}
    \leanok
    At each recorded source slot, apply the prescribed left and right
    endpoint operators at their actual owners. Other registers remain
    unchanged, and no source is prepared by these operations.
\end{definition}

\begin{theorem}[Locality of the endpoint operators]
    \label{thm:peps_actual_endpoint_locality}
    \lean{TNLean.PEPS.PairEffect.Word.eval_localEndpoint_tmul,
      TNLean.PEPS.PairEffect.SourceInventory.restrict_endpointWord_eval_eq_of_agree}
    \leanok
    \uses{def:peps_actual_endpoint_word}
    A single endpoint operation sends $x\otimes y$ to $Ax\otimes y$.
    If two endpoint families agree at all owners in a selected set of
    parties, their operator words restricted to that set are equal.
\end{theorem}

\begin{proof}\leanok
    The single endpoint action follows from the tensor product of its map
    with the identity. Induct on the source slots. A selected endpoint
    uses the same operator in both families; an unselected endpoint has
    no effect on the restriction. Composition and framing preserve these
    equalities.
\end{proof}

These results describe the actual local operators in
\cite[proof of Theorem~5.2]{OpenAI2026PolynomialPEPS},
\texttt{04-compression.tex}, lines 279--355 and 565--588.
